Low Earth orbit (LEO) satellite constellations are becoming a backbone of 6G
non-terrestrial networks, yet their open inter-satellite links (ISLs) are
exposed to eavesdropping, injection, and replay, while on-board routers offer
too little compute and memory to run conventional per-flow security such as
IPsec or DTLS. This paper presents 6FLAS, a lightweight, stateless security protection
mechanism that moves security state off the satellite and into the packet
itself. 6FLAS repurposes the 20-bit IPv6 Flow Label as a connectionless flow
authenticator and binds it to an SRv6 security segment, so each on-board router
verifies a single keyed message authentication code (MAC) over the routing
header using only a small per-direction replay window rather than a per-flow
security association table. The LEO security problem is formalized on a
resource-constrained directed graph; the 6FLAS packet format, per-hop
verification algorithm, and a security argument against eavesdropping,
forgery, man-in-the-middle, replay, and stateless-denial-of-service attacks are
provided. 6FLAS is evaluated with a reproducible analytical processing model, an
event-driven latency model of a 1584-satellite Walker-delta constellation,
and a Python software prototype that validates the structural ordering on
commodity hardware.
Analytically projected on a 1\,GHz space-grade core, 6FLAS reduces per-hop
security processing time by about $48\times$ versus IPsec ESP and $5.9\times$
versus an SRv6 HMAC baseline; the Python prototype confirms the structural
ordering on commodity hardware ($3.6\times$ over an SRv6-Sec baseline). 6FLAS
cuts per-packet byte overhead by $79\%$ and eliminates per-flow on-board state
entirely (zero bytes per flow versus $256$\,B for IPsec~ESP), while keeping the
end-to-end latency within propagation-dominated bounds. These results indicate
that packet-carried, topology-aware security is a practical fit for the
resource envelope of LEO on-board routers.
